\title{Diabetic Foot Ulcers Severity Grading from Colour Wound
Images Using Deep Learning and Tissue Mapped Explainable AI}

\author{Mustapha~Nasser%
\thanks{M. Nasser is with the Faculty of Science and Technology,
American International University-Bangladesh, Dhaka, Bangladesh
(student ID 211002042).}}

\maketitle

\begin{abstract}
No public diabetic foot ulcer dataset carries severity annotation, so
recent work derives severity labels from image content. This paper
implements one such derivation, clustering wound pixels by colour and
thresholding the necrotic fraction, and then subjects the resulting
labels to four independent validation checks rather than assuming they
are correct. All four fail. Agreement with two clinicians is
indistinguishable from chance ($\kappa \approx 0$), against an
inter-rater agreement of $\kappa = 0.25$ that bounds what any automated
method could achieve. A negative control returns a severe grade for
60.2\% of lesion-free skin images. We identify the mechanism: $k$-means
with $k=3$ cannot express the absence of a tissue type, so the darkest
region of any image is labelled necrotic by construction, including
shadow, wound depth and darker skin. To separate label quality from
pipeline quality, the identical architecture and protocol are applied to
an expert-annotated infection task, used strictly as a control and not as
a substitute for severity, where the same code reaches AUROC 0.81.
Severity grading from colour wound images is therefore limited by the
absence of clinical annotation rather than by model capacity, and we
release the validation protocol that establishes this.
\end{abstract}

\begin{IEEEkeywords}
Diabetic foot ulcer, label validation, data leakage, negative control,
wound image analysis, ordinal regression, explainable AI.
\end{IEEEkeywords}

\section{Introduction}

\IEEEPARstart{D}{iabetes} mellitus affects a large and growing share of the
world population, and the International Diabetes Federation reported 537
million adults living with the disease in 2021, with projections reaching
783 million by 2045 \cite{idf2021}. Among its complications, diabetic foot
ulcers (DFU) carry the heaviest clinical and economic burden. Between 25\%
and 34\% of diabetic patients develop an ulcer during their lifetime
\cite{armstrong2017}, and more than 85\% of diabetes-related lower-limb
amputations are preceded by a foot ulcer \cite{pecoraro1990}. In
lower-middle-income countries the cost of treatment can reach several times
a patient's annual income \cite{cavanagh2012}.

Clinicians grade DFU severity using systems such as the Wagner scale
\cite{wagner1981} or the University of Texas classification
\cite{armstrong1998}, both of which map the wound to an ordinal category
that determines the treatment pathway. Assessment is performed by visual
inspection supported by physical probing, which introduces inter-clinician
variability and is unavailable in settings without wound care specialists
\cite{vannetten2020}. Automating this assessment from a photograph would
extend triage capacity to primary care and remote settings.

Deep learning has been applied to DFU image analysis since 2018
\cite{goyal2018}, and progress accelerated after the Diabetic Foot Ulcer
Challenge series introduced standardised benchmarks \cite{cassidy2021}. A
2025 systematic review screening 4{,}653 articles found that machine
learning models achieve accuracies between 0.88 and 0.97 on binary
detection and on infection or ischaemia classification \cite{gomes2025}.
Severity grading, however, remains unsolved on public data. Studies that
attempt Wagner grading rely on private hospital datasets that cannot be
reproduced \cite{zhou2025,abebe2025}, or collapse adjacent grades because
labelled data is insufficient \cite{zhou2025}.

The obstacle is not model capacity but annotation. No publicly downloadable
DFU dataset carries severity labels. A natural response, and the one this
work initially pursued, is to derive severity from image content: segment
the wound into tissue types by colour, compute the proportion of each, and
apply threshold rules anchored to published grading criteria. The approach
is attractive because it requires no clinician time and scales to any
unlabelled corpus.

This paper reports what happens when that derivation is validated rather
than assumed correct. It does not recover severity, and we document both
the evidence and the mechanism. To establish that the failure lies in the
labels and not in the surrounding pipeline, the same architectures and
evaluation protocol are applied to a wound classification task that does
carry expert annotation, where they perform as expected. The conclusion is
therefore specific: severity grading from colour wound images is not
blocked by model capacity or image resolution, but by the absence of
clinically annotated ground truth.

A second finding emerged from the same experiments. During the study we
designed TGO-Net, a tissue-guided architecture that injects the derived
tissue proportions into a frozen convolutional representation through a
learned multiplicative gate whose scalar is initialised at zero, so the
network reports for itself how much of that signal it uses. On the derived
labels the gated model appears to perform extremely well, but we show this
is circular: the tissue proportions are the label-generating variable. On
expert labels, where no such circularity exists, the gate never stabilises
and the gated model performs slightly worse than a plain CNN. The
architecture is retained in this paper as a diagnostic instrument rather
than a performance contribution, because the behaviour of its gate is
itself evidence about the labels.

The contributions of this work are:

\begin{itemize}
\item A \textbf{colour-based tissue proportion labelling method} that
derives mild, moderate and severe labels from wound images without
ground-truth annotations, specified in full so that it can be reproduced
and independently assessed.

\item A \textbf{four-part validation protocol} for derived image labels,
combining expert agreement, a negative control on lesion-free images, a
stability test across augmented copies of the same source photograph, and a
corpus plausibility check. Each check fails independently, and each is
inexpensive relative to the cost of training on invalid labels.

\item \textbf{Identification of the mechanism} behind the failure:
$k$-means with $k=3$ partitions every image into three non-empty clusters
and therefore cannot express the absence of necrosis; the assumed luminance
ordering of tissue types is inverted; and all dark image regions, including
darker skin, are absorbed into the cluster labelled necrotic.

\item A \textbf{leakage-free evaluation protocol} for pre-augmented corpora,
using dihedral-invariant content hashing, photograph-level identity
recovery from the union of filename and pixel evidence, and patient-grouped
folds verified on the image table across five independent keys.

\item \textbf{TGO-Net}, comprising a Tissue Proportion Encoder, a
Cross-Modal Gate with a zero-initialised scalar, and a rank-consistent
CORAL head, presented with its parameter budget and with the gate scalar
used as an interpretable measure of modality reliance.

\item A \textbf{control experiment} applying the identical architectures,
preprocessing and evaluation protocol to an expert-annotated infection
task, isolating label provenance as the cause of failure rather than model
capacity, input resolution or training configuration.
\end{itemize}

The remainder of this paper is organised as follows. Section II reviews
related work. Section III states the critical gaps in prior methods.
Section IV lists the core contributions. Section V describes the
methodology. Section VI details the datasets and experimental setup.
Section VII reports results. Section VIII discusses them. Sections IX
through XI state limitations, future work and conclusions.

\section{Related Work}

\subsection{Severity Grading from Wound Images}

The most direct prior work on Wagner grading from images is Zhou et al.
\cite{zhou2025}, who applied Mask2Former to segment seven tissue types and
assign grades on 671 annotated images from a Chinese hospital, reaching AUC
0.94 on tissue segmentation. Grades 1 and 2 were collapsed and grades 0 and
5 excluded because of insufficient samples, and the dataset is not public.
Abebe et al. \cite{abebe2025} applied MobileNetV2 to full six-class Wagner
grading on 3{,}700 images from three Ethiopian hospitals, reporting 82\%
validation accuracy, again on private data and without cross-dataset
validation. Anilkumar et al. \cite{anilkumar2026} compared six
architectures on a 407-image South Indian clinical dataset labelled by
IWGDF criteria, with MobileNetV2 reaching 82\% accuracy and adjacent-grade
confusion identified as the dominant failure mode. Wang et al.
\cite{wang2024score} proposed ScoreDFUNet, classifying images into ulcer,
infection, normal and gangrene zones on 1{,}426 images with 95.34\%
accuracy, without addressing the full severity scale.

The pattern across these studies is consistent: severity grading works when
clinicians have graded the images, and every such dataset is private. None
reports a protocol for deriving severity labels from public unlabelled
data, and none validates such labels.

\subsection{Architectures for DFU Classification}

The standard architecture shifted from custom convolutional networks to
pretrained backbones after the challenge series demonstrated that
fine-tuned EfficientNet models outperform task-specific designs on
infection and ischaemia classification \cite{liu2022}. Subsequent work has largely
pursued attention mechanisms and feature fusion on the same benchmarks:
transformer-attention hybrids \cite{karthik2025}, attention-driven dense networks
, reconstruction residual networks with spatial-channel attention
\cite{wang2024farr}, handcrafted texture descriptors combined with learned embeddings
\cite{algaraawi2022}, and patch-based networks with context-preserving attention \cite{liu2021}.
Reported macro F1 on the DFUC-2021 infection and ischaemia task
clusters around 80\%.

These studies share a limitation relevant here: all evaluate on labels
supplied with the dataset, and none reports whether augmented copies of the
same photograph, or multiple wound sites from the same patient, were
separated across training and test splits.

\subsection{Explainability in DFU Models}

Explainability work in this domain applies SHAP, LIME and Grad-CAM to
binary detection, reporting high accuracies and heatmaps that localise
the ulcer region \cite{rathore2025,mahmud2025}. A recent
review found that only 23\% of DFU studies incorporate explainability
methods at all, and that none maps activation outputs to clinical wound
tissue features \cite{liu2025review}.

Across this literature, explainability outputs are generated at the
detection level and are not linked to validated tissue references. Our
analysis in Section VII-E addresses why doing so is harder than it appears
when the tissue reference is itself derived by clustering.

\subsection{Label Quality and Data Leakage in Medical Imaging}

Concern about label quality and evaluation integrity is not specific to
DFU. Reviews of machine learning in medical imaging repeatedly identify
optimistic evaluation arising from improper data partitioning, in
particular when multiple images originate from the same patient or the same
acquisition \cite{yap2021,silva2025}. Gomes et al. \cite{gomes2025} note
that only 67.3\% of DFU classification datasets are publicly accessible,
which limits independent verification of both labels and splits. Dhar et
al. \cite{dhar2024} report a pilot study on wound tissue segmentation and
observe that inter-annotator variation is substantial even among trained
annotators, which is consistent with the inter-rater agreement we measure
in Section VII-A.

To our knowledge no prior DFU study reports a negative control on
lesion-free images, or a stability analysis across augmented copies of the
same photograph, as a precondition for trusting derived labels.

\section{Critical Gaps and Limitations of Prior Methods}

Four gaps follow from the literature above, and the contributions in
Section IV map onto them one to one.

\textbf{Gap 1: derived labels are used without validation.} Where public
data lacks severity annotation, labels are derived from image content and
then used directly for training. No prior study reports whether the derived
labels agree with clinicians, whether they behave correctly on images
containing no lesion, or whether they are stable under transformations that
do not change the clinical picture. Contribution 2 supplies a validation
protocol; contribution 3 reports the mechanism when it fails.

\textbf{Gap 2: evaluation integrity on pre-augmented corpora is not
reported.} Public wound corpora are frequently distributed after
augmentation, so a split over image files does not separate photographs.
None of the reviewed studies states how augmented copies or per-patient
multiplicity were handled. Contribution 4 supplies a protocol that
addresses this and verifies it on the image table rather than on an
aggregated table.

\textbf{Gap 3: architectural additions are not tested against a
label-quality confound.} Where auxiliary features derived from the same
images are added to a network, an improvement may reflect the auxiliary
feature carrying label information rather than clinical signal.
Contribution 5 introduces an architecture whose gate scalar makes modality
reliance directly readable, and contribution 6 supplies the control
condition needed to interpret it.

\textbf{Gap 4: no control separates label quality from pipeline quality.}
When a method underperforms, the literature does not distinguish between
inadequate labels and an inadequate model. Contribution 6 applies the
identical pipeline to an expert-annotated task on comparable images,
isolating the variable.

\section{Core Contributions}

The contributions listed in Section~I are developed as follows: the
labelling method in Section~V-B, the validation protocol in Section~V-C
with results in Section~VII-A, the mechanism in Section~VII-B, the
evaluation protocol in Section~V-D, TGO-Net in Sections~V-E and~V-F, and
the control experiment in Section~VII-C.